\section{Why the disciples can look up}

In the Gospel of the First Sunday of Advent, the nations can hear the sea. People are already faint with fear of what is still to come. Jesus turns to his disciples and tells them to raise their heads when these things begin. They too can hear the sea. What they have heard from him lets them look for someone in the midst of the terror. The Son of Man is coming, and Jesus calls his approach their redemption.

Before the Gospel the singers ask God to show them mercy. In the Gospel the disciples are to see the Son of Man in majesty. At the altar the priest gives thanks to the Father for the divine majesty that the Son and Spirit share with him. A Christian who has made the singers' petition his own will want to know how he may look forward to the sight of the Judge. He has just asked, in the Collect, for rescue from the danger of his own sins.

Gregory the Great asks his hearers to consider whom they love when they imagine losing the world they know. John Chrysostom finds a reason for hope in Paul's choice of a word. Paul tells people who have let themselves fall asleep that their \emph{salvation} is near. Their first believing lies behind them, and some of their first eagerness may have gone with it. Paul still speaks to them as to runners within reach of the prize.

The coming of Christmas raises a different question. Paul writes to Christians for whom the Lord's birth has already taken place. How can the Church use his words to await that birth? Thomas Aquinas observes that people who expect a feast can speak with those who waited for what the feast celebrates. He hears in the Advent lesson the voices of generations who believed Christ would come. There is a second future tense at Communion. The earth \emph{shall} yield its fruit. Robert Bellarmine reads those words of the psalm as a promise that Mary will bear a son. Beside the Communion antiphon, Prosper Guéranger asks Christians to recognize this child in the Lord they have received. Someone still looking ahead to the feast can give thanks for a gift already made.

\begin{maptable}
Introit & Ps 24(25):1--4 & A man threatened with ridicule entrusts himself to God and asks to learn his paths.\\
Collect & \textit{Excita} (Stir up) & The Church calls on the Lord to come and free her from the dangers her members' sins have brought upon them.\\
Epistle & Rom 13:11--14a & Paul wakes those who first believed, tells them what to put off, and bids them put on the Lord Jesus Christ.\\
Gradual & Ps 24(25):3--4 & The promise to all who wait returns with the request to learn a path.\\
Alleluia & Ps 84:8 (85:7) & The people ask the Lord to let them see his mercy and to give them his salvation.\\
Gospel & Lk 21:25--33 & People faint for fear. The disciples may lift their heads. The fruit on the trees tells them summer is near.\\
Offertory & Ps 24(25):1--3 & The Introit antiphon returns while the gifts are offered.\\
Secret & \textit{Haec sacra} (These holy gifts) & The priest asks for the gifts to cleanse those who draw near to their author.\\
Preface & Most Holy Trinity & The priest and people lift their hearts and give thanks. Father, Son and Spirit receive equal adoration.\\
Communion & Ps 84:13 (85:12) & God will give goodness and the earth will yield fruit.\\
Postcommunion & \textit{Suscipiamus} (May we receive) & The people ask for mercy as they prepare to celebrate their restoration.\\
\end{maptable}

Psalm numbers follow the Missal. Parentheses give the corresponding modern psalm number and, where it differs, the modern verse. The places of the biblical speakers, the episodes recalled and the accounts of when the passages were written appear in Scriptural Date and Location (Appendix~\ref{sec:date-location}).
